\section{Исследуемые наноструктуры}

В данной работе были исследованы структуры, состоящие из кремния на подложке из $SiO_2$. Характерные параметры системы приведены в таблице \ref{tbl:params}. Схематично система изображена на рисунке \ref{fig:scheme}. Голубой цвет отвечает подложке из $SiO_2$, а красный цвет -- кремниевым наноструктурам.

\begin{table}[H]
	\centering
	\begin{tabular}{|c|c|}
		\hline
		Параметр & Значение\\
		\hline
		\hline
		Длина волны источника & 1.0 мкм--2.0 мкм \\
		\hline
		Длина волновода, $L$ & 10 мкм \\
		\hline
		Ширина волновода, $w$ & 0.6 мкм \\
		\hline
		Высота волновода, $h$ & 0.25 мкм, 0.4 мкм \\
		\hline
		Высота нанодиска, $h$ & 0.25 мкм, 0.4 мкм \\
		\hline
		Радиус нанодиска, $r$ & 0.1 мкм--1.0 мкм\\
		\hline
		Расстояние между нанодиском и волноводом, $d$ & $-$0.09 мкм--1.0 мкм\\
		\hline
		Количество нанодисков & 1--20 \\
		\hline
		Расстояние между соседними нанодисками & 0 мкм--0.4 мкм \\
		\hline
	\end{tabular}
	\caption{Характерные параметры исследуемой системы}
	\label{tbl:params}
\end{table}

\begin{figure}[h]

	\begin{subfigure}[b]{\textwidth}
		\centering
		\includegraphics[width=\textwidth]{img/scheme_xy3}
		\caption{Проекция XY}
		\label{fig:scheme_xy}
	\end{subfigure}
	\begin{subfigure}[b]{\textwidth}
		\centering
		\includegraphics[width=\textwidth]{img/scheme_xz}
		\caption{Проекция XZ}
		\label{fig:scheme_xz}
	\end{subfigure}
	\begin{subfigure}[b]{\textwidth}
		\centering
		\includegraphics[width=.7\textwidth]{img/scheme_yz3}
		\caption{Проекция YZ}
		\label{fig:scheme_yz}
	\end{subfigure}
	\caption{Схематичное изображение исследуемой системы}
	\label{fig:scheme}

\end{figure}